\documentclass[UTF8, a4paper, 12pt]{ctexart}

%% ============================================
%% 中文经济学论文 LaTeX 模板
%% 适配国内《经济研究》《管理世界》等期刊格式
%% ============================================

\usepackage[top=2.54cm,bottom=2.54cm,left=3.17cm,right=3.17cm]{geometry}
\usepackage{setspace}
\usepackage{indentfirst}
\setlength{\parindent}{2em}
\usepackage{microtype}

% 数学
\usepackage{amsmath, amssymb}
\numberwithin{equation}{section}

% 表格与图表
\usepackage{booktabs}
\usepackage{graphicx}
\usepackage{longtable}
\usepackage{array}
\usepackage{float}
\usepackage{subcaption}

% 交叉引用
\usepackage{hyperref}
\hypersetup{colorlinks=true, linkcolor=black, citecolor=black}

% 自定义命令
\newcommand{\E}{\mathbb{E}}
\newcommand{\Var}{\mathrm{Var}}
\newcommand{\Cov}{\mathrm{Cov}}
\DeclareMathOperator{\plim}{plim}
\DeclareMathOperator{\argmax}{argmax}
\DeclareMathOperator{\argmin}{argmin}

% 作者信息
\title{\textbf{中文标题}\thanks{基金项目支持（项目编号：XXXX）。感谢匿名审稿人的宝贵意见，文责自负。）}}
\author{
张三$^{1}$\footnote{通讯作者：zhangsan@edu.cn} \quad 李四$^{2}$ \\
$^{1}$\textit{北京大学经济学院} \\
$^{2}$\textit{清华大学经济管理学院}
}
\date{}

\begin{document}
\maketitle
\thispagestyle{empty}

% ============================================
% 摘要
% ============================================
\begin{center}
\textbf{摘\quad 要}
\end{center}
\vspace{1ex}

本文研究......利用......数据，采用......方法，实证检验了......。基准回归结果显示，......。异质性分析表明，......。机制检验发现，......。研究结论表明，......

\medskip
\noindent
\textbf{关键词：} 研究主题；核心变量；研究方法；政策含义 \\
\textbf{JEL分类号：} O10; C50; R10

\newpage

% ============================================
% 1. 引言
% ============================================
\section{引言}

\subsection{研究背景与问题}
[介绍研究背景，引出研究问题]

\subsection{研究意义}
[从理论和实践两个维度说明研究意义]

\subsection{研究创新}
[列举本文的主要贡献：]
\begin{itemize}
    \item 可能的研究贡献一：从......角度分析......
    \item 可能的研究贡献二：使用......数据进行实证检验......
    \item 可能的研究贡献三：揭示了......机制
\end{itemize}

\subsection{文章结构}
[说明后续各部分的安排]

% ============================================
% 2. 文献综述与理论分析
% ============================================
\section{文献综述与理论分析}

\subsection{核心概念界定}
[明确本文涉及的核心概念]

\subsection{理论框架}
[基于相关理论构建分析框架]

\subsection{文献回顾与评述}

\subsubsection{关于......的研究}
[回顾相关实证研究，指出其发现和局限]

\subsubsection{关于......的研究}
[回顾相关实证研究]

\subsection{研究假说}
[基于理论和文献提出可检验的研究假说：]

\textbf{假说H1：} ......

\textbf{假说H2：} ......

% ============================================
% 3. 数据与研究设计
% ============================================
\section{数据与研究设计}

\subsection{数据来源}
[说明数据来源、时间范围、样本选择]

\subsection{变量定义}
[清晰定义被解释变量、解释变量、控制变量]

\begin{table}[H]
\caption{变量定义表}
\label{tab:var_def}
\centering
\begin{tabular}{cll}
\toprule
变量类型 & 变量名称 & 定义说明 \\
\midrule
被解释变量 & $Y$ & [定义] \\
\midrule
核心解释变量 & $X$ & [定义] \\
\midrule
控制变量 & $Z_1$ & [定义] \\
 & $Z_2$ & [定义] \\
 & $Z_3$ & [定义] \\
\bottomrule
\end{tabular}
\end{table}

\subsection{描述性统计}
[报告主要变量的描述性统计量]

\begin{table}[H]
\caption{描述性统计}
\label{tab:desc_stat}
\centering
\begin{tabular}{lccccc}
\toprule
变量 & 观测值 & 均值 & 标准差 & 最小值 & 最大值 \\
\midrule
$Y$ & 1,234 & 0.567 & 0.234 & 0 & 1 \\
$X$ & 1,234 & 4.321 & 1.567 & 0 & 10 \\
$Z_1$ & 1,234 & 2.345 & 0.876 & 1 & 5 \\
$Z_2$ & 1,234 & 0.678 & 0.465 & 0 & 1 \\
$Z_3$ & 1,234 & 0.789 & 0.408 & 0 & 1 \\
\bottomrule
\end{tabular}
\end{table}

\subsection{模型设定}
[设定回归模型]

本文基准回归模型设定如下：

\begin{equation}
\label{eq:baseline}
Y_{it} = \alpha + \beta X_{it} + \gamma Z_{it} + \mu_i + \lambda_t + \varepsilon_{it}
\end{equation}

其中，$i$ 表示地区，$t$ 表示年份。$\mu_i$ 和 $\lambda_t$ 分别为地区和年份固定效应，用于控制不可观察的地区和时间异质性。

% ============================================
% 4. 实证结果
% ============================================
\section{实证结果}

\subsection{基准回归结果}
[报告基准回归结果]

\begin{table}[H]
\caption{基准回归结果}
\label{tab:baseline}
\centering
\begin{tabular}{lccc}
\toprule
 & \multicolumn{3}{c}{被解释变量 $Y$} \\
\cmidrule(lr){2-4}
 & (1) & (2) & (3) \\
\midrule
核心解释变量 $X$ & 0.123^{***} & 0.098^{**} & 0.087^{*} \\
 & (0.032) & (0.041) & (0.045) \\
\midrule
控制变量 & 否 & 是 & 是 \\
地区固定效应 & 否 & 是 & 是 \\
年份固定效应 & 否 & 否 & 是 \\
\midrule
观测值 & 1,234 & 1,234 & 1,234 \\
$R^2$ & 0.123 & 0.245 & 0.389 \\
\bottomrule
\end{tabular}
\begin{minipage}{\linewidth}
\smallskip
\textit{注：} 括号内为稳健标准误，$^{***}$、$^{**}$、$^{*}$分别代表1\%、5\%、10\%的显著性水平。
\end{minipage}
\end{table}

\subsection{稳健性检验}

\subsubsection{替换核心变量}
[使用替代变量进行检验]

\subsubsection{子样本分析}
[按不同子样本进行分析]

\begin{table}[H]
\caption{异质性分析结果}
\label{tab:heterogeneity}
\centering
\begin{tabular}{lcc}
\toprule
 & \multicolumn{2}{c}{分组样本} \\
\cmidrule(lr){2-3}
 & 高$Z$组 & 低$Z$组 \\
\midrule
核心解释变量 $X$ & 0.156^{**} & 0.045 \\
 & (0.062) & (0.038) \\
\midrule
观测值 & 617 & 617 \\
$R^2$ & 0.345 & 0.298 \\
\midrule
组间系数差异检验 & \multicolumn{2}{c}{$p$值 = 0.023} \\
\bottomrule
\end{tabular}
\end{table}

\subsection{机制检验}
[检验可能的作用机制]

% ============================================
% 5. 结论与启示
% ============================================
\section{结论与启示}

\subsection{主要结论}
[总结本文的主要发现]

\subsection{政策启示}
[基于研究发现提出政策建议]

\subsection{研究局限}
[承认研究的局限性]

\subsection{未来研究方向}
[提出未来可拓展的方向]

% ============================================
% 参考文献
% ============================================
\newpage
\begin{thebibliography}{99}
\addtolength{\itemsep}{-1ex}

bibitem{} 张三, 李四. 论文标题[J]. 期刊名称, 2023, 38(1): 1-15.

bibitem{} Author, A. and Author, B. Title of the paper[J]. Journal Name, 2023, 45(2): 123-145.

bibitem{} World Bank. World Development Report 2023[M]. Washington D.C., 2023.

\end{thebibliography}

% ============================================
% 附录
% ============================================
\newpage
\appendix
\section*{附录}

\subsection*{表A1：变量描述性统计（续）}
% [完整的描述性统计]

\subsection*{表A2：稳健性检验——替换核心变量}
% [额外的稳健性检验]

\end{document}